\label{ch:info_als_energie}
% Source document: 255

\epigraph{Information is topologically frozen kinetic energy.}{Doc.\,255}

\section{Two Types of Energy in \texorpdfstring{$\tilde{T} \cdot m = 1$}{T·m = 1}}

The fundamental relation of the FFGFT $\tilde{T} \cdot m = 1$ enforces a natural
division of energy: static energy as topologically enclosed
kinematics and free kinetic energy as motion through $T^4$.

\textbf{Static energy:} Energy bound to rest mass $E_0 = mc^2$
corresponds to topologically stable winding numbers on $T^4$. It is discrete, does
not contribute to thermodynamic entropy, and remains conserved as long as the
topology remains intact.

\textbf{Dynamic kinetic energy:} Free kinetic energy --- for
massless particles the only form of energy ($E = hf$), for massive
particles the part beyond the rest mass. It is continuous,
thermodynamic and entropy-producing.

\section{Information as Frozen Kinematics}

From this division follows without additional postulates:

\begin{keyresult}
Information is topologically frozen kinetic energy. A
winding configuration $(n_\theta, n_\varphi)$ on $T^4$ is a
kinetic energy that is secured against dissipation by the topology of the torus.
It cannot be lost without the topology breaking ---
and topological transitions are quantized in FFGFT.
\end{keyresult}

Consequence: Landauer's principle describes the price for the
\emph{thawing} of frozen kinematics. The erasure of a bit is the
enforced transfer of a topologically secured configuration into
another --- the phase space content released in the process goes to the bath.

\section{Quantization of Information}

The discreteness of the winding numbers enforces the quantization of
information. There is no half bit --- the winding number is an integer.
That is not a postulate, but a direct consequence of the torus topology.

The winding quantum $\Delta w = 1$ is thus the fundamental
information unit of the FFGFT --- smaller is not possible, because the topology
knows no non-integer winding numbers.

\section{M.\,M.~Vopson from an FFGFT Perspective}

Vopson's equivalence \textbf{[V1]} $E_{\text{info}} = mc^2_{\text{Bit}}$ can be
derived from FFGFT --- but only for the special case of a specific scale.
The energy of a winding quantum on scale $L$ is $\hbar c / L$, and the
corresponding equivalent mass is $m = \hbar c / (Lc^2) = \hbar / (Lc)$.
That is not a universal bit mass --- it is scale-sensitive.

Vopson's \textbf{[V1]} result is the limiting case $L \to L_P$ (Planck scale):
then $m_{\text{bit}} \to m_P$, the Planck mass. Matzke's derivation
$m_{\text{bit}} = m_P \ln 2 / (2\pi)$ is the Landauer special case at
Planck temperature. Both are limiting cases of the FFGFT formula, not
universal laws.

\section{The Casimir Effect as Evidence}

The Casimir effect shows that the vacuum itself carries information:
The vacuum energy density at the characteristic $\xi$ length $L_\xi$
matches the Casimir energy density. That is no coincidence ---
it is the FFGFT prediction that $\xi$ determines both the geometric scale
of the torus windings and the vacuum energy scale (Doc.\,025).